\scanpage{45}\mspage{51}
\unsure{changt|changement} de base pour
\diagram{\begin{tikzcd}[ampersand replacement=\&]
 Z'\arrow[r,"p"]\arrow[d,"f'"'] \& Z\rlap{$\;F$}\arrow[d,"f"]\\
 S\arrow[r,"\mathrm{id}"]\&S
\end{tikzcd}\quad $.\quad\big]$}

\uline{Remarque} : On retrouve le fait que $i^{\unsure{*}}i_*(F_Y)\xrightarrow{\sim}F_Y$ est iso, pour $i$ une imm. fermée. \deleted{\unsure{Retrouvons}} De plus avec les notations de Dér 4 (g), \deleted{\unsure{si on}} \deleted{\unsure{met}} $j^*(i_*(F_Y))$ se calcule par le th. de chgt de base (cor. 2 ci-dessus), qui s'applique dans la situation commutative
\diagram{\begin{tikzcd}[ampersand replacement=\&]
 X'\arrow[r,"q"]\arrow[d,"f'"']\arrow[dr,phantom,"\mathrm{cart}" description] \& X\arrow[d,"f"]\\
 Y'\arrow[r,"p"']\&Y
\end{tikzcd}}
quand $p$ est une imm. ouverte (et $\unsure{\gamma|r}=\mathrm{id}$) (alors $X'_{/y'}\to X_{/p(y')}$ est \unsure{$=$} un iso !), et on trouve, \unsure{revenant|venant} aux notations de Der 4 (g)
\diagram{\begin{tikzcd}[ampersand replacement=\&]
 \varnothing\arrow[r,"\unsure{q|g}"]\arrow[d,"f'"']\& Y\arrow[d,hook,"i"]\\
 U\arrow[r,hook,"j"']\& X
\end{tikzcd}}
ici : $j^*(i_{\deleted{\unsure{!|*}}*}\unsure{(F_Y)|[F_Y]}\simeq f'_*(\varphi)$, où $\varphi=\unsure{g|q}^*(F_Y)$ est l'``unique'' élément de la catégorie (équivalente à une cat. ponctuelle, par Der 1) $\mathbb D(\varnothing)$. D'autre part $f'_*$ transforme objet final en objet final, d'où $j^*i_*(F)$ est objet final.

Inversement, soit $F$ dans $\mathbb D(X)$, telle que $j^*(F)=e_{\mathbb D(U)}$. Considérons
\[
 F\xrightarrow{u}i_*i^*F\qquad(\mathrm{adj})
\]
je dis que c'est un \uline{iso}. En effet, par l'axiome
% archive sheet number: 44
